% ── Appendix ───────────────────────────────────────────────────────────────
\appendices
\section{Author Contributions}

The project was divided into three pipeline stages, each owned by one team
member. Table~\ref{tab:contributions} summarises the primary files and
responsibilities; narrative details follow.

\begin{table}[h]
\centering
\caption{Team member ownership by pipeline stage.}
\label{tab:contributions}
\renewcommand{\arraystretch}{1.25}
\begin{tabular}{@{} l l @{}}
\toprule
\textbf{Member} & \textbf{Stage} \\
\midrule
My Nguyen Tra       & Data \& Preprocessing \\
Long Trinh Le Hoang & Stage 1: CNN Training \\
Hao Nguyen Thien    & Stage 2: Sequence Model \& Deployment \\
\bottomrule
\end{tabular}
\end{table}

\subsection*{My Nguyen Tra — Data \& Preprocessing}

My Nguyen Tra was responsible for everything from raw DICOM files to
model-ready tensors. Her tasks included:
\begin{itemize}
    \item Extraction of CQ500 archives and verification of DICOM file integrity.
    \item Label generation from \texttt{reads.csv} via majority vote
          ($\geq$\,2/3 readers) and broadcast of study-level labels to
          every slice.
    \item Stratified train/val/test splitting (70/15/15) across 491
          studies, preserving subtype proportions in each split.
    \item HU preprocessing pipeline: \texttt{RescaleSlope}/\texttt{RescaleIntercept}
          conversion, clipping to [$-$1000, 3000]\,HU, brain-window
          stacking of three adjacent slices $(s-1,\,s,\,s+1)$ as a
          3-channel float32 image.
    \item Augmentation pipeline (\texttt{HFlip}, \texttt{ShiftScaleRotate},
          \texttt{RandomErase}, \texttt{RandomCrop} for training;
          centre crop for validation).
    \item Study-based \texttt{ICHStudyDataset} with centre-weighted
          slice sampling and class-imbalance analysis
          (\texttt{pos\_weight} per class).
\end{itemize}

\subsection*{Long Trinh Le Hoang — Stage 1: CNN Training}

Long Trinh Le Hoang owned the CNN backbone, training loop, scheduler,
and evaluation framework. His tasks included:
\begin{itemize}
    \item Implementation of \texttt{DenseNet121ICH} and
          \texttt{DenseNet169ICH} backbones with an
          \texttt{AdaptiveAvgPool2d(1)} head and a \texttt{get\_features()}
          method for downstream feature extraction.
    \item \texttt{BCEWithLogitsLoss} with per-class \texttt{pos\_weight}
          as the primary loss; Adam optimiser
          (lr\,=\,5$\times 10^{-4}$, weight\_decay\,=\,2$\times 10^{-5}$).
    \item WarmRestart cosine-annealing scheduler (SGDR):
          10-epoch warmup plateau followed by \texttt{T\_mult=2} cycles.
    \item Mixed-precision training (AMP), gradient clipping, MLflow
          experiment tracking, early stopping, and checkpoint saving.
    \item Per-class AUC, macro-AUC, sensitivity and specificity metrics
          with bootstrap 95\,\% confidence intervals.
\end{itemize}

\subsection*{Hao Nguyen Thien — Stage 2: Sequence Model \& Deployment}

Hao Nguyen Thien owned the BiGRU sequence model and the Docker-based
inference API. His tasks included:
\begin{itemize}
    \item Feature extraction script to convert Stage-1 logits into
          per-study \texttt{.pt} files consumed by Stage 2.
    \item Two-sub-network BiGRU architecture: SM1 (FC\,+\,BiGRU) and
          SM2 (dilated Conv1D\,+\,BiGRU with skip connection), summed
          to produce refined per-slice logits.
    \item Stage-2 training with Adam\,+\,CosineAnnealingLR for 40 epochs.
    \item Grad-CAM saliency visualisation pipeline
          (\texttt{reports/figures/gradcam.png}).
    \item FastAPI \texttt{/api/v1/predict} endpoint accepting DICOM
          uploads and returning per-class probabilities.
    \item Multi-stage Docker build (\texttt{make serve-docker}) enabling
          full inference with no local Python installation.
\end{itemize}
